\paragraph{纤维同伦判别、最大非可缩纤维与 Lee--Yang 三次数域链}

\begin{theorem}[纤维奇偶与同伦型的等价判别]\label{thm:killo-fold-fiber-homotopy-parity-equivalence}
\leanpartial{paper\_killo\_fold\_fiber\_homotopy\_parity\_equivalence}{the requested Lean signature only supplies one-way implications `badModThreeComponent -> contractibleCase` and `allComponentsAvoidBadModThree -> sphereCase` inside `FiberIndependenceComplexClassificationData`, so the converse directions in the stated parity equivalences are not derivable in the current repository}
设 \(m\ge 1\)，\(x\in X_m\) 为稳定类型。若其候选位点诱导图分解为路径的离散并
\[
\Gamma_x \cong \bigsqcup_{i=1}^{r} P_{\ell_i},
\qquad
d_m(x)=\prod_{i=1}^{r} F_{\ell_i+2},
\]
其中 \(F_n\) 为 Fibonacci 数。定义纤维独立集复形
\[
K_x:=\Ind(\Gamma_x),\qquad
\tau(x):=\sum_{i=1}^{r}\left\lfloor \frac{\ell_i+1}{3}\right\rfloor .
\]
则有：
\begin{enumerate}
  \item \(|K_x|\) 可缩，当且仅当 \(d_m(x)\) 为偶数；
  \item \(|K_x|\simeq S^{\tau(x)-1}\)，当且仅当 \(d_m(x)\) 为奇数。
\end{enumerate}
\end{theorem}
\begin{proof}
对路径并分解的独立集复形，存在如下同伦二分：若存在某个分量满足
\(\ell_i\equiv 1\pmod 3\)，则 \(|K_x|\) 可缩；若所有分量都满足
\(\ell_i\not\equiv 1\pmod 3\)，则 \(|K_x|\simeq S^{\tau(x)-1}\)。

另一方面，\(F_n\) 在模 \(2\) 下的 Pisano 周期为 \(3\)，故
\[
F_n\equiv 0\pmod 2
\iff
3\mid n.
\]
代入 \(n=\ell_i+2\) 得
\[
F_{\ell_i+2}\ \text{为偶}
\iff
\ell_i\equiv 1\pmod 3.
\]
因此
\[
d_m(x)\ \text{为奇}
\iff
\forall i,\ \ell_i\not\equiv 1\pmod 3,
\]
以及
\[
d_m(x)\ \text{为偶}
\iff
\exists i,\ \ell_i\equiv 1\pmod 3.
\]
将该奇偶判别与同伦二分逐条拼合即得结论。
\end{proof}

\begin{theorem}[最大非可缩纤维的模 \(6\) 相图判别]\label{thm:killo-fold-max-noncontractible-fiber-mod6}
定义
\[
D_m:=\max_{x\in X_m} d_m(x),
\qquad
\widetilde D_m:=\max\{d_m(x):x\in X_m,\ |K_x|\not\simeq *\}.
\]
若前三层纤维谱满足
\[
D_{2k}=F_{k+2},\qquad D_{2k+1}=2F_{k+1},
\]
\[
D^{(2)}_{2k}=F_{k+2}-F_{k-4}\ (k\ge 6),\qquad
D^{(2)}_{2k+1}=2F_{k+1}-F_{k-4}\ (k\ge 5),
\]
\[
D^{(3)}_{2k}=F_{k+2}-F_{k-3}\ (k\ge 6),\qquad
D^{(3)}_{2k+1}=2F_{k+1}-F_{k-4}-F_{k-8}\ (k\ge 8),
\]
则对所有 \(m\ge 17\)，有完整分类
\[
\widetilde D_m=
\begin{cases}
D_m, & m\equiv 0,4 \pmod 6,\\[2mm]
D_m^{(2)}, & m\equiv 1,5 \pmod 6,\\[2mm]
D_m^{(3)}, & m\equiv 2,3 \pmod 6.
\end{cases}
\]
\leanverified{paper\_killo\_fold\_max\_noncontractible\_fiber\_mod6}
\end{theorem}
\begin{proof}
由定理 \ref{thm:killo-fold-fiber-homotopy-parity-equivalence}，\(|K_x|\not\simeq *\) 当且仅当 \(d_m(x)\) 为奇数，故
\[
\widetilde D_m=\max\{d_m(x):x\in X_m,\ d_m(x)\ \text{为奇}\}.
\]

先看偶数层 \(m=2k\)。由 \(D_{2k}=F_{k+2}\)，当且仅当 \(3\nmid(k+2)\) 时 \(D_{2k}\) 为奇，因此该情形下 \(\widetilde D_{2k}=D_{2k}\)。若 \(3\mid(k+2)\)（即 \(k\equiv 1\pmod 3\)），则 \(D_{2k}\) 为偶。再看
\[
D^{(2)}_{2k}=F_{k+2}-F_{k-4},
\]
此时 \(k+2\equiv k-4\equiv 0\pmod 3\)，两项均偶，故 \(D^{(2)}_{2k}\) 仍偶。进一步
\[
D^{(3)}_{2k}=F_{k+2}-F_{k-3},
\]
其中 \(k-3\equiv 1\pmod 3\)，故 \(F_{k-3}\) 为奇，从而偶减奇为奇，得 \(\widetilde D_{2k}=D^{(3)}_{2k}\)。将 \(k\equiv 0,1,2\pmod 3\) 映射到 \(m=2k\equiv 0,2,4\pmod 6\)，得到偶数层判别。

再看奇数层 \(m=2k+1\)（\(m\ge 17\Rightarrow k\ge 8\)，故第三层闭式可用）。有
\[
D_{2k+1}=2F_{k+1},
\]
恒为偶数，故最大层不可实现 \(\widetilde D\)。考察
\[
D^{(2)}_{2k+1}=2F_{k+1}-F_{k-4}.
\]
若 \(3\nmid(k-4)\)，则 \(F_{k-4}\) 为奇，故偶减奇为奇，得 \(\widetilde D_{2k+1}=D^{(2)}_{2k+1}\)。若 \(3\mid(k-4)\)（即 \(k\equiv 1\pmod 3\)，等价于 \(m\equiv 3\pmod 6\)），则 \(D^{(2)}_{2k+1}\) 为偶。此时 \(k-8\equiv 2\pmod 3\)，故 \(F_{k-8}\) 为奇，进而
\[
D^{(3)}_{2k+1}=2F_{k+1}-F_{k-4}-F_{k-8}
\]
为偶减偶减奇，故为奇，得 \(\widetilde D_{2k+1}=D^{(3)}_{2k+1}\)。与 \(m\equiv 1,3,5\pmod 6\) 对应后即得奇数层判别。
\end{proof}

\begin{theorem}[Lee--Yang 分支点平方根系数的三次数域刻画]\label{thm:killo-leyang-kappa-square-cubic-field}
考虑谱四次多项式
\[
\Pi(\lambda,y)=\lambda^4-\lambda^3-(2y+1)\lambda^2+\lambda+y(y+1).
\]
设其非平凡 Lee--Yang 分支点 \((\lambda_\star,y_\star)\) 满足
\[
16\lambda_\star^3-9\lambda_\star^2+1=0,
\]
并定义
\[
u:=\kappa_\star^2
:=
-2\frac{\partial_y\Pi}{\partial_{\lambda\lambda}^2\Pi}\Big|_{(\lambda_\star,y_\star)}
=\frac{2(3\lambda_\star^2-1)}{3(\lambda_\star^2-1)}.
\]
则 \(u\) 满足不可约三次方程
\[
324u^3-540u^2+333u-74=0,
\]
其判别式为
\[
\Disc\!\bigl(324u^3-540u^2+333u-74\bigr)
=-2^{6}\cdot 3^{9}\cdot 37.
\]
特别地，该三次多项式的 Galois 群为 \(S_3\)。
\end{theorem}
\begin{proof}
由
\[
u=\frac{2(3\lambda^2-1)}{3(\lambda^2-1)}
\]
解得
\[
\lambda^2=\frac{3u-2}{3(u-2)}.
\]
又由分支方程 \(16\lambda^3-9\lambda^2+1=0\)（且分支点处 \(\lambda\neq 0\)）除以 \(\lambda^2\) 得
\[
\lambda=\frac{9\lambda^2-1}{16\lambda^2}.
\]
代入上式可化为
\[
\lambda=\frac{3(2u-1)}{4(3u-2)}.
\]
再将该表达代回 \(16\lambda^3-9\lambda^2+1=0\)，清分母即得
\[
324u^3-540u^2+333u-74=0.
\]
有理根检验排除一次因子，故该三次在 \(\QQ\) 上不可约。判别式按三次公式计算即为
\(-2^{6}\cdot 3^{9}\cdot 37\)，其非平方，因此不可约三次的 Galois 群必为 \(S_3\)。
\end{proof}

\begin{theorem}[Lee--Yang 振幅常数 \(B^2\) 的最小多项式与有理互构]\label{thm:killo-leyang-B-square-cubic-and-rational-intertwine}
在定理 \ref{thm:killo-leyang-kappa-square-cubic-field} 的分支点处，记
\[
c^2=\kappa_\star^2,\qquad
B:=\frac{c\sqrt{y_\star}}{\lambda_\star},\qquad
|B|\approx 0.9634999692,
\]
并令 \(b:=B^2\)。则：
\begin{enumerate}
  \item \(b\) 满足不可约三次方程
  \[
  162b^3+1593b^2+1998b+1184=0;
  \]
  \item 其判别式为
  \[
  \Disc\!\bigl(162b^3+1593b^2+1998b+1184\bigr)
  =-2^{2}\cdot 3^{9}\cdot 11^{2}\cdot 37\cdot 109^{2},
  \]
  因而该三次的 Galois 群为 \(S_3\)；
  \item 若 \(u=\kappa_\star^2\)，则
  \[
  b
  =
  -2\,
  \frac{24708u^2-28628u+7733}{51084u^2-44412u+8735},
  \]
  并且
  \[
  \QQ(b)=\QQ(u)=\QQ(\lambda_\star).
  \]
\end{enumerate}
\end{theorem}
\begin{proof}
由分支条件 \(\partial_\lambda\Pi(\lambda,y)=0\) 可解得
\[
y=\frac{4\lambda^3-3\lambda^2-2\lambda+1}{4\lambda},
\]
并与 \(\Pi(\lambda,y)=0\) 联立回到分支三次
\(
16\lambda^3-9\lambda^2+1=0
\)。
由定义
\[
b=B^2=\frac{c^2y}{\lambda^2}
=\frac{\kappa_\star^2\,y}{\lambda^2}
=\frac{2(3\lambda^2-1)}{3(\lambda^2-1)}\cdot\frac{y}{\lambda^2},
\]
可知 \(b\) 为 \(\lambda\) 的有理函数。将该有理表达记作 \(b=p(\lambda)/q(\lambda)\)，对方程组
\[
16\lambda^3-9\lambda^2+1=0,\qquad q(\lambda)\,b-p(\lambda)=0
\]
对 \(\lambda\) 消元（例如使用 Sylvester 结式）即得
\[
162b^3+1593b^2+1998b+1184=0.
\]
有理根检验排除一次因子，故其在 \(\QQ\) 上不可约。判别式按三次公式计算得到
\(
-2^{2}\cdot 3^{9}\cdot 11^{2}\cdot 37\cdot 109^{2}
\)，非平方，故 Galois 群为 \(S_3\)。

再由定理 \ref{thm:killo-leyang-kappa-square-cubic-field} 中
\[
\lambda=\frac{3(2u-1)}{4(3u-2)},
\]
代回 \(b=p(\lambda)/q(\lambda)\) 并约分，得到显式有理关系
\[
b=-2\,
\frac{24708u^2-28628u+7733}{51084u^2-44412u+8735}.
\]
因此 \(b\in\QQ(u)\)。另一方面，\(u\) 与 \(b\) 的最小多项式都为不可约三次，故
\[
[\QQ(u):\QQ]=[\QQ(b):\QQ]=3.
\]
由包含关系立即得到 \(\QQ(b)=\QQ(u)\)。又因 \(u\in\QQ(\lambda_\star)\) 且 \(\lambda_\star\in\QQ(u)\)，遂有
\(\QQ(b)=\QQ(u)=\QQ(\lambda_\star)\)。
\end{proof}
